% ============================================================
\section{Limitações}\label{sec:limitacoes}
% ============================================================

As limitações abaixo são as conhecidas até o planejamento e a etapa \etapa{corpus}. Cada uma diz
o que pode distorcer os resultados e o que foi feito para reduzir o efeito. A lista é
atualizada a cada etapa.

\begin{description}
    \item[O tema é só um indicador aproximado do sentido.] Na P4, o tema do artigo faz o papel
    do sentido de uma palavra-alvo: \emph{banco} em Computação é, em geral, banco de dados, mas
    um artigo de Computação pode falar do Banco Central. Uma separação por tema pode, então,
    subestimar a separação por sentido (ocorrências com o ``sentido errado'' para o tema) ou
    refletir diferenças de tema que não são de sentido (vocabulário vizinho). Mitigações: a
    anotação manual de cerca de 20 ocorrências por alvo (\cod{senses.csv}), analisada ao lado do
    tema, e a comparação com os controles, que têm a mesma alocação por tema mas um sentido
    estável: o que os controles mostrarem de ``separação por tema'' é o efeito do tema sem
    mudança de sentido.

    \item[O Qwen3 fragmenta o português.] Só 28\% dos vértices do estudo de viabilidade são
    palavras inteiras (Seção~\ref{sec:viabilidade}). Com isso, boa parte dos vértices é
    fragmento sem significado próprio, cuja vizinhança lexical reflete a ortografia e não o
    sentido, e as palavras-alvo ficaram restritas às que são um token só (12 de 16; saíram
    \emph{corrente}, \emph{onda}, \emph{célula} e \emph{núcleo}). As conclusões sobre
    fragmentos e palavras inteiras são reportadas separadamente, e as interpretações
    semânticas se restringem a palavras inteiras de conteúdo. A generalização para modelos com
    tokenizers melhores para o português (Tucano, GPT-2 português) não é testada.

    \item[Predefinições removidas deixam frases incompletas.] A primeira frase de muitos
    artigos usa predefinições para pronúncia, etimologia, datas e nomes em outras línguas
    (``X (em grego: \ldots) é\ldots''). Elas são descartadas na limpeza, e a limpeza tira os
    restos visíveis (parênteses vazios, vírgulas repetidas), mas algumas frases ficam mais
    curtas ou com construções pouco naturais. O modelo vê esse texto como contexto, então parte
    das primeiras posições dos parágrafos iniciais tem um contexto um pouco diferente do texto
    original. O efeito deve ser pequeno (só atinge algumas frases iniciais), e a faixa de
    posição entra como controle nas análises.

    \item[A busca por categorias é truncada.] A árvore de categorias da Wikipédia em português
    não é uma taxonomia limpa: em poucos níveis, uma categoria de ciência chega a pessoas,
    lugares e obras. Por isso a busca vai só até a profundidade 2 e para em 600 títulos por
    tema. O corte não é aleatório. A API lista os membros de cada categoria na ordem da chave de
    ordenação (em geral, alfabética), e a busca em largura abre as subcategorias nessa ordem,
    nível por nível; quando o limite é atingido, as subcategorias que ainda estão na fila nunca
    são abertas. O embaralhamento com semente (Seção~\ref{sec:corpus}) vem depois e só escolhe
    os 150 artigos de cada tema entre esses 600 títulos. Na execução principal, os oito temas
    atingiram o limite, e a busca abriu só uma parte das subcategorias que encontrou: de 7 de
    79 (Computação) a 98 de 439 (Música); em Política, 23 de 443. Em Biologia, Política e
    Computação, o corte veio antes de terminar a profundidade 1, e nenhum título veio da
    profundidade 2; em Computação, 494 dos 600 títulos são membros diretos da raiz. Em Música e
    Futebol, a busca parou no meio de uma lista de subcategorias por país (em ``Música da
    Austrália'' e ``Futebol das Bahamas''), de modo que, dessa lista, só os países do começo do
    alfabeto foram abertos. (Números obtidos refazendo a busca a partir do cache da etapa
    \etapa{corpus}, sem rede; a busca refeita reproduz a categoria e a profundidade de todos os
    artigos examinados.) Ficaram de fora subáreas inteiras: em Biologia, 17 das 33
    subcategorias diretas da raiz, como Reprodução, Simbiose e Teorias biológicas; em Política,
    33 de 55, como Movimentos sociais e Políticas públicas; em Computação, 43 de 49, como
    Inteligência artificial, Hardware e Mineração de dados (uma delas, Redes de computadores,
    entra como categoria extra). Entre os 150 artigos mantidos de
    Futebol, 61 vêm da profundidade 2: 32 de categorias por continente e 29 de países de
    Afeganistão a Azerbaijão; não há nenhum de Brasil ou Portugal. A profundidade dos artigos
    mantidos também varia de um tema para outro: em Computação, 123 dos 150 são membros diretos
    da raiz; em Geografia, 105 dos 150 vêm da profundidade 2. Os temas representam, portanto,
    as subcategorias mais próximas da raiz e alfabeticamente primeiras, e não a área inteira.
    Como o objetivo é ter temas bem separados, e não estimar o vocabulário de cada área, isso é
    aceitável, mas os nomes dos temas devem ser lidos como ``artigos próximos da categoria
    raiz''. A comparação entre representações não é afetada, porque os vértices são as mesmas
    ocorrências em todas as redes; as análises que usam o tema como rótulo (P3, e P4 com o tema
    como indicador do sentido) herdam o recorte. O mesmo corte limita o
    descarte de páginas multitema: um título só é reconhecido como multitema se estiver na
    lista truncada de mais de um tema, e um artigo que também pertence à parte não aberta de
    outra árvore passa como se fosse de um tema só. O descarte de páginas multitema e de
    biografias reduz a mistura, ao custo de excluir páginas de fronteira entre temas.

    \item[Precisão \emph{bfloat16}.] O modelo roda em bf16 (8 bits de mantissa), e os vetores
    são guardados em bf16. O cosseno em \emph{float64} é exato \emph{para os vetores
    guardados}, mas esses vetores já carregam o arredondamento das contas do modelo: duas
    execuções com sequências de comprimentos diferentes podem dar estados ligeiramente
    diferentes para o mesmo prefixo (Seção~\ref{sec:mod-sequencias}). Vizinhanças decididas por
    diferenças de cosseno menores que esse ruído são, na prática, empates. Mitigações: a cópia
    por grupo de prefixo (que torna exatos os empates que deveriam ser exatos), o diagnóstico da
    maior diferença dentro dos grupos (\difMaxPrefixo) e o teste de sensibilidade com
    $\eps=10^{-4}$. Na GPU da execução (RTX A5500, 24 GiB), o modelo caberia em
    \emph{float32}, mas isso não foi adotado: o bf16 é o tipo nativo do modelo, os vetores
    continuam guardados em bf16, e o passo de arredondamento do bf16 nas magnitudes do bloco 36
    já é da ordem da diferença medida dentro dos grupos de prefixo.

    \item[A Wikipédia provavelmente está no pré-treino.] O Qwen3 foi pré-treinado em cerca de 36
    trilhões de tokens de fontes não detalhadas~\cite{qwen3}, e a Wikipédia em português quase
    certamente faz parte delas. O modelo pode, então, ter visto (e em parte memorizado) os
    parágrafos analisados. Isso não invalida a análise da geometria das representações, mas
    pode torná-las mais ``confiantes'' do que seriam em texto novo (por exemplo, uma previsão
    do próximo token mais acertada e, no bloco 36, uma organização mais forte pelo próximo
    token). Não há como controlar isso sem texto certamente fora do pré-treino; fica registrado
    como ressalva de interpretação.

    \item[Um modelo, uma língua, uma amostra.] Os resultados são de um único modelo
    (Qwen3-4B-Base, numa revisão fixada), de texto enciclopédico em português e de uma amostra
    desenhada (as frequências na amostra refletem as cotas, e não o texto natural). As faixas
    de frequência no corpus e a análise da P3 só no núcleo reduzem o efeito do desenho, mas
    não o eliminam.
\end{description}
